% Author interpretation written after the completed, preregistered reporting
% pass. The report's immutable generated snippets and all raw records remain
% unchanged. Numerical bindings: INTERPRETATION_BINDINGS120.json.
\subsection{Economic gains and the strong conventional comparison}
\label{sec:economicinterpretation120}

The gains in Table~\ref{tab:economic120} concern a concrete investment
decision. Investment increases the two next-period stocks by $a/2$ and $a/4$,
whereas the continuous innovation shifts them by $(z,-z)$. The innovation
therefore changes their imbalance without changing their aggregate. The
decision trades the convex current adjustment cost against stock tracking,
the low-stock penalty, the capacity constraint and the original risk-sensitive
continuation. At a common initial state, the current state cost cancels from
the comparison. It is the action-dependent continuation difference, together
with the adjustment-cost difference, that must justify an intervention.
These are normalized original-objective units. The experiment supplies no
monetary or consumption-equivalent welfare calibration.

Numerical displays in this discussion summarize retained exact rational
intervals; the exact endpoints determine every sign and price condition.
All seven learned representations improve on their own completed incumbents
at both risk levels under each of L, P and D. The independent mean-gain lower
bounds are strictly positive in all 42 such cells. For ReLU, P and L return
the same policies: the mean original-policy gains are approximately
$[212.675,213.885]\times10^{-6}$ at $\alpha=1/2$ and
$[162.411,163.501]\times10^{-6}$ at $\alpha=3/4$.
The corresponding independent mean-regret upper bounds fall from
$215.906$ to $2.626$ and from $165.361$ to $2.406$, respectively, in the
same $10^{-6}$ units. The original-action RMS-error upper bounds fall from
$0.01470$ to $0.001621$ and from $0.01286$ to $0.001551$.
These are bounds on error relative to an optimal original action; an exact
optimal action was not supplied as an observed benchmark.

The magnitude of an improvement also depends on the initial actor's quality.
The spline gains are larger than the ReLU gains because the completed spline
incumbent has a larger original loss in this catalogue. They do not establish
a superior final spline policy. More decisively, the online semi-analytic
incumbent has a strictly lower sample mean original objective than every
learned completed incumbent in all 14 comparisons. Every refinement leaves
that online policy unchanged on the 64 states at both risks. Its independent
mean-regret upper bounds are already $1.898\times10^{-6}$ and
$1.851\times10^{-6}$. Online also retains a strictly lower original objective
than ReLU under L and P: the independently resolved online advantages are
$[0.1506,1.2864]\times10^{-6}$ and
$[0.06473,1.0451]\times10^{-6}$ at the two risks. ReLU with D reaches
intervals overlapping the online policy, without resolving a strict mean
advantage. All 512 ReLU--hinge policy pairs are exactly identical, and both
representations pay the complete learned-feature acquisition charge. This
control rules out a neural-specific policy advantage within the executed
equal-feature comparison.

\subsection{What the new comparison changes}
\label{sec:incrementalevidence120}

The comparison with inherited improvement is more demanding than the
comparison with an unrefined learned incumbent. P changes 391 of its 1024
returned actions and accepts 815 strict improvement steps. L changes 388
actions and accepts 807 steps. Their final actions differ on only 12 of the
1024 common state/representation/risk cases. Nine of the 16 cell-level
policy comparisons are exact identities; all seven remaining independent
mean-value intervals contain zero. The retained intervals include both
directions. Thus this experiment resolves gains over learned incumbents
while leaving the incremental mean economic advantage of P over L
unresolved.

The action-scaled bound does affect an actual stopping decision. At stored
state index 29, $\alpha=1/2$, the Poly3 candidate is accepted with $m=4$:
the new original-policy difference upper bound is approximately
$-1.05403\times10^{-5}$, whereas the separately valid convex-only scalar
upper bound is $7.53783\times10^{-6}$. This is the same candidate in P and
D, recorded twice; it is one economic example, not two independent
observations. The mixed upper is active at 5961 of 9016 paid precision
levels, but is necessary for acceptance at only those two recorded levels.
The transferred incumbent tail-defect bound is inactive in all 2290 acquired
density pairs. Its theorem remains valid; no observed reduction in incumbent
tail work is attributed to it. These diagnostics establish use of the new
comparison at a specified paid resolution. They do not establish failure of
a traditional method after a further refinement or measure a counterfactual
time saving.

The complete costs in Table~\ref{tab:cost120} impose a further qualification.
P costs more than L in every one of the 16 primary cells, with ratios from
$1.143$ to $1.821$. It is also slower in all 48 matched fresh clocks; their
complete-cost ratios range from $1.048$ to $2.040$. All three clocks, including
their different magnitudes, remain reported. D is faster than P in every
primary cell and all 48 matched clocks. Against L, D is faster in 41 clocks
and slower in seven, with 14 favorable and two adverse cell medians. Its
independent original-policy mean is strictly better than L and P in five
cells---ReLU and hinge at both risks, and spline at $\alpha=3/4$---with two
exact online identities and nine unresolved comparisons. The conventional
direct proposal is therefore an essential comparator to the new
three-proposal rule.

Work counts and elapsed cost rank these procedures differently. Summed over
the same 16 cells, L purchases 2,032,754 compiled terminal rows, P purchases
1,758,340, and D purchases 14,608,426. The corresponding sums of complete
standalone 64-state costs are 94.478, 132.729 and 80.302 seconds. P retains
2257 rejected candidate attempts, including 785 exhausted precision ladders;
D retains 633 rejections, including 194 exhausted ladders. Identical-action
shortcuts remain in these counts and are distinguished from purchased
precision levels. The price of a point-information interface cannot be
inferred from its point count alone. These measured results retain separate
compiler, proposal, risk-query and certificate counters and the actual
cohort clocks; they do not combine overlapping counters into a universal
information unit or identify an unmeasured causal timing decomposition.

\subsection{Population scope and actual adoption conditions}
\label{sec:adoptioninterpretation120}

The frozen population analysis supplies a positive mean-gain lower bound for
all 42 learned-incumbent comparisons, with the simultaneous event-and-mean
family coverage stated in Table~\ref{tab:population120}. The event-staircase
lower bounds range from $2.661\times10^{-6}$ to $14.192\times10^{-6}$.
The ordinary bounded-mean lower bounds alone are zero; the predeclared
positive-gain events supply the positive lower conclusions. The upper
intervals remain broad. The six online comparisons have observed zero gains
and lower bound zero; the finite sample does not prove that refinement is
unnecessary everywhere. The inference concerns the declared finite-uniform
initial-state distribution conditional on frozen actors and the ideal
independent entropy assumption. It does not turn the 992 same-sample
cross-representation contrasts into a new population comparison family.

A gain over a weak incumbent can pay for refinement without making its
representation the preferred method. Using the actual 64-state workload
and its own acquisition-plus-cohort cost, P improves the cost-augmented
ReLU incumbent mean for computing prices in
$[0,0.001902379\ldots)$ at $\alpha=1/2$ and
$[0,0.001769390\ldots)$ at $\alpha=3/4$, in original-objective units per
second. The exact rational open endpoints, not their printed decimal
approximations, define these sets. Every one of the 32 ReLU comparisons
against the full conventional menu, across the four predeclared amortization
volumes and after intersection across all three clocks, has an empty
sufficient-price set. The corresponding intersections of upper
``not ruled out'' conditions are empty as well. Across all representations,
the 17 nonempty sufficient menu-price regions belong to completed
conventional policies. These are conditions on the recorded sample and the
declared cost scenario, not evidence of population adoption or measured
throughput scaling. When a cheaper policy has a worse original objective,
time savings can compensate only in the cost-augmented calculation; they do
not improve the unpriced original Bellman objective.

The deployment comparison includes actual production certificates, refused
proposals, acquisition, loading and cohort bookkeeping. The 64 primary
cohort clocks sum to 331.550 seconds; the complete primary process, including
three timing subprocesses and its other work, takes 841.861 seconds.
The separate independent rational reconstruction takes 1806.153 seconds
and is paid as research validation rather than silently assigned zero or
amortized away in the deployment-price formula. Source/artifact checks,
imports, process startup and final receipt work remain in the retained
process and launcher ledger; they are not claimed as separately measured
per-arm cohort components. Acquisition-plus-cohort sums charge standalone
acquisition to each hypothetical deployment and therefore differ from
the actual shared acquisition process. The mathematical progress is a valid
action-dependent comparison and a certified rule that changes decisions
under the original objective. The completed evidence identifies the
remaining computational task precisely: preserve that validity while
reducing paid decision work relative to strong conventional proposals and
resolving advantages beyond ordinary equal-feature representations.
